\section{Impact sur ma pratique future}

Bien que la connaissance factuelle et la rigueur soient des fondements essentiels de la pratique médicale, compte tenu des enjeux vitaux, il serait réducteur de s'arrêter là. Ne risquons-nous pas de perdre l'essence de la relation humaine si nous échouons à décoder ce qui se transmet au-delà du verbal ? Le langage non-verbal est, en effet, aussi un vecteur fondamental de la communication médecin-patient. À cet égard, la célèbre théorie d'Albert Mehrabian \parencite{demetersante2022}, qui attribue une part prépondérante au non-verbal ($55$\% pour le langage corporel) dans la communication des attitudes et des émotions, confirme l'impératif de développer cette compétence. Apprendre à interpréter les signaux corporels du patient devient alors une composante nécéssaire à notre formation. Ce besoin est amplifié face aux patients dont l'expression est limitée par des facteurs tels que l'âge, la langue, la culture ou la pathologie. La danse, en tant qu'expression corporelle par excellence, représente ainsi une éducation précieuse à l'observation fine de l'autre.\\

\todoH{FloraStyle}{red!30}{\ passage trop stigmatisant. MG médecine environnementale !!! apporter des nuances }{\mbox{Au-delà de l'émotion artistique}}, cette expérience invite à explorer le concept du One Health défini par une approche systémique reconnaissant que la santé humaine, la santé animale et l'équilibre des écosystèmes sont fondamentalement interdépendants. L’antibiorésistance illustre parfaitement la nécessité de cette approche. Face à cette menace, les mesures traditionnelles de la médecine humaine (contrôle des prescriptions, recherche, hygiène hospitalière) restent insuffisantes. Le One Health impose une vision systémique et multidisciplinaire. Concrètement, cela implique d'optimiser les prescriptions médicales et vétérinaires, mais aussi d'améliorer le traitement des \todoH{FloraStyle}{red!30}{\ perturbateurs endocriniens dans les eaux usées pas que AB --> il y a plein d'exemples possibles mais questions de nombre de mots, je ne vais pas en parler}{eaux usées} pour y réduire la concentration d’antibiotiques et de bactéries résistantes. L'éducation et la prévention auprès de tous ces professionnels (soignants, vétérinaires, agriculteurs) sont également primordiales. Heureusement, des transitions s'opèrent déjà dans l'enseignement professionnel : en Belgique, l'utilisation des huiles essentielles \parencite{siep_formation_524} en agronomie est enseignée comme alternative aux antibiotiques en élevage, avec des plans d'actions aux résultats encourageants selon \textcite{celagri_antibioresistance}.\\
Je suis convaincue que l’intégration du One Health représente l’avenir d'une médecine engagée et préventive, vers laquelle notre formation tend de plus en plus. À notre échelle, des gestes cliniques modestes (une anamnèse approfondie, un examen rigoureux et une prescription réfléchie avant tout examen complémentaire) contribuent déjà à une pratique plus durable et à la transformation globale de notre manière de soigner.\\


Certes, cette vision peut sembler idéaliste, mais elle est porteuse d’espoir. Même si cette transition vers une médecine plus responsable ne se fera pas en un jour, elle se fera. Ce "micro-rappel" précieux nous recentre sur l’essentiel et nous ouvre à une réflexion plus large sur notre mission de soignant.